% ============================================================
% 第4章 特征工程
% 【负责人：数据线（4.1 基础事件特征已预填）+ 建模线（4.2–4.4 补充）】
% 赛题要求“特征工程”为说明文档必备内容。
% ============================================================
\section{特征工程}

\subsection{基础事件特征（F0，已完成）}

基于清洗后的事件表，对 24 类事件编码按“一车一窗口”生成
\textbf{24 类 $\times$ 3 形态} 的基础特征矩阵，三形态对应三种计数口径
（见 3.3 节）：\texttt{raw\_count}（原始报警条数）、\texttt{segments30}
（30 秒并段后的段数）、\texttt{days}（发生天数）。特征命名遵循
\texttt{evt\_<编码>\_<形态>\_<窗口>} 规范（如 \texttt{evt\_11401\_count\_20d}），
改名必须升级 \texttt{feature\_version}。

特征口径要点：
\begin{itemize}
  \item 训练窗特征取自 6/1--6/21（20 天），推理窗特征取自 7/11--7/31
        （最近 20 天），\textbf{两窗口同定义}——不得把训练时的 20 天计数
        换成 60 天计数；
  \item 无该类记录计 0，同时附覆盖标记（\texttt{quality\_flags}：
        \texttt{imu\_not\_observed} / \texttt{not\_observed}），
        未知不误作 0；
  \item 每个特征列的定义、窗口与缺失规则登记于《已生成特征字典》
        （节选见附录 B），全量字典随代码交付。
\end{itemize}

在基础特征上聚合出 9 个基线特征组（\texttt{fatigue}、\texttt{distraction}、
\texttt{lane}、\texttt{collision\_warning}、\texttt{overspeed}、
\texttt{blindspot}、\texttt{device}、\texttt{history}、
\texttt{event\_observed\_days}），用于逻辑回归基线（见第 6 章），五折折外
AUC 0.7766（95\% CI 0.726--0.826），说明清洗后的基础特征已具备
可用的事故区分能力。

\subsection{轨迹暴露量特征（F1，待 2.0 轨迹数据）}

\begin{TODO}
负责人：数据线。计划内容：真实公里数、驾驶小时（\texttt{sum(run\_time)/3600}，
增量口径，先按车排序、去精确重复、处理乱序与跨分片，并用相邻时间差交叉验证）、
观测覆盖天数、事件率 = 事件计数 / 暴露量。硬规则：暴露量分母必须与事件计数
来自同一时间范围；分母无效时率置 null 并附原因，不得用 0 或极小常数掩盖。
当前 2.0 轨迹数据 403 阻塞，结构验证基于 1.0 分片。
\end{TODO}

\subsection{IMU 派生特征（F2，待 2.0 IMU 数据）}

\begin{TODO}
负责人：数据线 + 建模线。计划内容：IMU 覆盖率与质量门（不输出物理动作结论，
直至单位与采样机制确认）、有交集行驶时段内的动作率（急刹/急加速段统计，
分母为轨迹 $\cap$ IMU 覆盖时段）、设备差异标准化。硬规则：缺失 IMU 不得填为
安全；不得因覆盖不全剔除目标车辆；v4 死区窗内不提取动作特征、缺口日特征置
缺失（见 3.4/3.6 节）。
\end{TODO}

\subsection{特征降维、聚类与选择}

\begin{TODO}
负责人：建模线。按 2026-09-18 会议共识撰写：
(1) 特征融合——画像数据与事件数据融合，通过聚类合并相似指标
（如夜间行驶时长与里程）降低维度；
(2) 维度控制——样本仅 500 台，特征维度控制在 10--20 个左右防过拟合；
(3) 特征选择比较规则——比较特征时固定模型，比较模型时固定特征；
全部选择只在训练折内完成，不得用全体标签预筛特征后再声称独立交叉验证。
达标后补入实验记录与结论。
\end{TODO}
